% ==========================================================================
% 第二章  燃煤机组汽轮机系统状态监测建模方法（v4）
% 对齐杨昊东式章节功能与真实VHP工程数据；本文件为Ch2审计事实源。
% ==========================================================================

\section{燃煤机组汽轮机系统状态监测建模方法}
\label{sec:gru_monitoring_v4}

\subsection{引言}

燃煤机组汽轮机系统在宽负荷、频繁变工况条件下运行时，设备压力、温度和流量等状态参数不仅受到设备健康状态影响，还会随入口蒸汽边界、负荷水平及操作调节持续变化。若直接将某一设计值或固定历史均值作为正常状态基准，容易把正常工况迁移误判为设备异常。因此，建立能够随运行边界动态变化的健康状态参考，是后续利用状态偏差信息开展故障诊断的基础。

本文按照汽轮机实际热力过程将系统划分为VHP、一次再热、HP、二次再热、IP、LP及凝汽器等设备级建模对象，并统一采用“操作变量$U$+热力边界$B\rightarrow$设备状态$Y$”的建模方式。模型输入仅使用当前设备能够直接获得的真实DCS操作量和边界量，不将上游模型预测输出继续串联为下游模型输入，从而避免多设备串联预测造成误差累积。考虑汽轮机过程参数存在明显的时间迟延与历史依赖，本章采用门控循环单元（Gated Recurrent Unit，GRU）建立设备级健康状态预测模型，并以VHP作为代表性对象开展完整工程算例。

本章首先说明汽轮机系统设备级建模对象及其输入输出边界；随后给出GRU动态建模方法及训练策略；在此基础上构建VHP健康状态模型，并通过2024年5月正常运行数据训练、2024年6月独立测试，与ANN、TCN和iTransformer-small进行对比，最后分析不同状态参数的可建模性及模型适用边界。

\subsection{汽轮机系统设备级状态监测建模对象}
\label{subsec:turbine_device_objects_v4}

汽轮机热力系统主通流可概括为“主蒸汽--VHP--一次再热--HP--二次再热--IP--LP--凝汽器”。各设备之间通过蒸汽介质持续传递热力状态，同时还与抽汽回热、给水和冷端系统形成耦合。为了使设备健康模型具有明确物理边界，本文不采用整机统一黑箱输入输出，而是依据介质流向和实际DCS测点对系统进行设备级划分。

设第$m$个设备在时刻$t$的操作变量、热力边界和状态输出分别为$\mathbf U_m(t)$、$\mathbf B_m(t)$和$\mathbf Y_m(t)$，则设备健康状态映射可表示为
\begin{equation}
\hat{\mathbf Y}_m(t)=f_m\left(\mathbf U_m(t-L+1:t),\mathbf B_m(t-L+1:t);\theta_m\right),
\end{equation}
式中：$L$为历史时间窗口长度；$\theta_m$为模型参数；$\hat{\mathbf Y}_m(t)$为给定当前操作和边界条件下的设备健康状态预测值。模型仅用于学习健康运行数据中的条件响应关系，不在本章直接承担故障分类任务。

设备边界的确定遵循三个原则：一是输入量应位于设备上游或外部边界，避免使用本级近出口量作为输入形成信息泄漏；二是输出量应能够反映设备本体及直接热力接口状态；三是所有输入输出均对应实际可获取DCS测点。对于相邻设备，上一设备的真实DCS出口测量量可以作为下一设备的实际入口边界，但不使用上一模型的预测值进行级联。

依据上述原则，本文当前构建的汽轮机系统状态监测对象包括VHP、RH1、HP、RH2、IP、LP和两个凝汽器等，主要对象及边界见表~\ref{tab:system_objects_v4}。

\begin{table}[htbp]
    \centering
    \caption{汽轮机系统主要设备级状态监测对象}
    \label{tab:system_objects_v4}
    \begin{tabular}{p{1.3cm}p{2.6cm}p{4.4cm}p{4.4cm}p{1.5cm}}
        \toprule
        编号 & 对象 & 主要输入边界 & 主要状态输出 & 数据可用性 \\
        \midrule
        M1 & VHP & 主蒸汽P/T、阀位、负荷、主汽流量 & VHP本体、一级抽汽、冷再接口及缸体热状态 & 完整 \\
        M2 & RH1 & VHP/冷再入口、再热调节相关边界 & 一次高温再热出口P/T及相关状态 & 可建模 \\
        M3 & HP & 一次高温再热入口P/T及工况边界 & HP本体及后续热力接口状态 & 完整 \\
        M4 & RH2 & HP排汽及二次再热入口边界 & 二次高温再热出口P/T及相关状态 & 可建模 \\
        M5 & IP & 二次高温再热入口及工况边界 & IP本体及抽汽相关状态 & 可建模 \\
        M6 & LP & IP/连通管相关入口及冷端边界 & LP本体及低压抽汽状态 & 可建模 \\
        M7/M8 & CND1/2 & LP排汽、循环水及冷端边界 & 真空、凝结水和冷端相关状态 & 可建模 \\
        \bottomrule
    \end{tabular}
\end{table}

上述设备划分用于统一全文“设备级健康状态建模”的物理接口。考虑本研究当前已完成并具有完整独立测试结果的对象为VHP，后续算法对比与状态监测结果分析均以VHP为代表性算例，其余设备用于说明该方法可按相同边界规则扩展，而不在本章宣称已经完成同等独立实验验证。

\subsection{GRU动态健康状态建模方法}
\label{subsec:gru_method_v4}

\subsubsection{GRU基本结构}

循环神经网络能够通过隐藏状态在相邻时间步之间传递历史信息，但传统RNN在较长时间序列训练中容易出现梯度消失或梯度爆炸。GRU通过更新门和重置门控制历史信息保留与当前输入写入，在较少参数规模下能够学习工业过程中的短期动态和一定历史记忆\cite{Cho2014GRU}。

对于时刻$t$的输入向量$\mathbf x_t$和上一时刻隐藏状态$\mathbf h_{t-1}$，GRU更新门$\mathbf z_t$和重置门$\mathbf r_t$分别定义为
\begin{equation}
\mathbf z_t=\sigma(W_z\mathbf x_t+U_z\mathbf h_{t-1}+\mathbf b_z),
\end{equation}
\begin{equation}
\mathbf r_t=\sigma(W_r\mathbf x_t+U_r\mathbf h_{t-1}+\mathbf b_r),
\end{equation}
式中：$\sigma(\cdot)$为Sigmoid激活函数；$W_z$、$U_z$、$W_r$、$U_r$为可学习权重；$\mathbf b_z$和$\mathbf b_r$为偏置项。更新门决定上一时刻隐藏状态有多少被保留，重置门控制历史信息在候选状态计算中的参与程度。

候选隐藏状态为
\begin{equation}
\tilde{\mathbf h}_t=\tanh\left(W_h\mathbf x_t+U_h(\mathbf r_t\odot\mathbf h_{t-1})+\mathbf b_h\right),
\end{equation}
最终隐藏状态更新为
\begin{equation}
\mathbf h_t=(1-\mathbf z_t)\odot\mathbf h_{t-1}+\mathbf z_t\odot\tilde{\mathbf h}_t,
\end{equation}
式中$\odot$表示逐元素乘法。通过门控机制，GRU能够在负荷变化、阀位调节以及蒸汽热力边界变化过程中保留对当前状态预测有价值的历史信息。

本文采用固定长度历史窗口构造设备输入序列。对于VHP，窗口长度为24个采样点，即24 min。设单时刻输入维数为$p$，则一个样本的输入矩阵为
\begin{equation}
X_t=[\mathbf x_{t-L+1},\ldots,\mathbf x_t]^{\mathrm T}\in\mathbb R^{L\times p}.
\end{equation}
GRU对整个窗口进行编码，取末时刻隐藏状态作为当前时刻的动态特征，并通过线性层同步预测$q$个设备状态量。

\subsubsection{多输出回归与损失函数}

汽轮机设备状态量同时包含压力、温度等不同物理量，原始量纲和数值范围差异较大。为避免大数值量纲在损失中占据主导，训练阶段首先使用训练集统计量对输入和输出分别进行标准化。对于第$j$个输出变量，标准化形式为
\begin{equation}
y_j^{*}=\frac{y_j-\mu_j}{s_j},
\end{equation}
式中$\mu_j$和$s_j$分别为训练集第$j$个变量的均值与标准差。验证集和测试集均只使用训练集得到的$\mu_j$和$s_j$，避免未来数据参与统计量估计。

模型采用标准化空间中的平均均方误差作为训练目标：
\begin{equation}
\mathcal{L}_{\mathrm{MSE}}=\frac{1}{N_bq}\sum_{i=1}^{N_b}\sum_{j=1}^{q}\left(y_{ij}^{*}-\hat{y}_{ij}^{*}\right)^2,
\end{equation}
其中$N_b$为批样本数，$q$为输出维度。标准化后，各输出维度在训练阶段具有相近数值尺度，使多输出任务能够在统一损失中共同参与优化。

在GRU输出后，本文引入LayerNorm对单样本隐藏特征进行归一化。LayerNorm根据同一层神经元特征计算均值和方差，其训练与测试阶段计算方式一致，并适用于循环网络隐藏状态稳定化\cite{Ba2016LayerNorm}。随后采用Dropout(0.05)进行轻度正则化，在训练过程中随机抑制部分隐藏特征，降低特征之间过度共适应的风险\cite{Srivastava2014Dropout}。由于本文网络规模较小，Dropout率设置较低，目的在于保持主要动态信息的同时提供必要正则化，而非通过较强随机失活改变网络表达能力。

\subsubsection{BPTT与参数优化}

GRU训练过程采用时间反向传播（BPTT）。在前向阶段，同一组GRU参数在$L$个时间步上重复使用，输入序列依次更新隐藏状态并形成多输出预测；计算回归损失后，梯度从输出层返回隐藏状态，并沿时间展开结构反向累积至各时间步共享参数，实现历史信息路径和当前输入路径的联合优化。

本文采用AdamW优化器进行参数更新。AdamW将权重衰减与自适应梯度更新过程解耦，可在使用Adam类自适应学习率的同时保持独立的权重衰减机制\cite{Loshchilov2019AdamW}。此外，训练中采用梯度范数裁剪限制异常梯度放大，并利用ReduceLROnPlateau在验证损失长期不改善时降低学习率，使模型在训练后期以更小步长搜索稳定参数区域。上述训练策略均服务于小型设备级网络的稳定优化，不作为本文独立算法创新点。

\subsection{基于GRU的汽轮机设备级状态监测模型构建}
\label{subsec:device_model_construction_v4}

在明确汽轮机设备边界及GRU动态建模原理后，本节进一步给出面向系统级应用的统一模型构建方法。该方法的重点不在于为不同设备设计完全不同网络，而在于依据真实热力流程确定各设备的边界输入和状态输出，并以一致的数据处理和训练协议建立可比较的健康状态参考。

\subsubsection{设备级输入输出定义}

对于任一设备模型，输入序列可统一写为
\begin{equation}
\mathbf x_t=[\mathbf U_t,\mathbf B_t],
\end{equation}
其中$\mathbf U_t$表示负荷、阀位、调节量等操作变量，$\mathbf B_t$表示设备入口蒸汽状态及必要的外部边界。输出向量$\mathbf y_t$由设备本体状态及直接热力接口状态组成。对于相邻设备，上一设备真实DCS出口测点可以进入下一设备的$\mathbf B_t$，但上一模型预测值不参与下一模型输入。

该结构的作用在于把正常工况变化解释为“给定$U+B$条件下$Y$应达到的健康响应”。后续若实际$Y$与健康预测$\hat Y$发生持续偏离，则偏差信息可进一步用于构造故障诊断节点的动态特征，而无需在状态监测阶段直接设定故障阈值。

\subsubsection{数据处理与样本构造}

设备原始DCS数据首先进行有效值筛选，并剔除计划停机、检修和明确无有效采集的时间区间。对正常运行数据按自然时间顺序构造长度为$L$的连续滑动窗口，窗口内部不允许跨越停机、无效数据或大时间断点。

输入和输出标准化参数仅由训练数据计算。对于VHP月度独立测试实验，2024年5月数据用于模型训练及内部验证，2024年6月数据作为完全独立测试集，不参与网络权重更新、标准化参数估计和模型选择。通过这种划分方式，测试结果反映模型从一个月份推广至下一独立月份的状态参考能力。

\subsubsection{网络结构与训练协议}

本文设备级GRU统一采用单层GRU编码历史窗口，隐藏维度为32，随后依次连接LayerNorm、Dropout和线性输出层。以VHP为例，实际训练网络为
\begin{equation}
24\times16\rightarrow \mathrm{GRU}(32)\rightarrow \mathrm{LayerNorm}\rightarrow \mathrm{Dropout}(0.05)\rightarrow \mathrm{Linear}(27),
\end{equation}
可训练参数总数为5755。这里Linear(27)对应训练阶段27个候选状态输出；完成状态点筛选后，后续统计分析和故障诊断映射统一使用24个状态点，因此不能将实际训练网络改写为Linear(24)。

优化器采用AdamW，初始学习率为$1\times10^{-3}$，权重衰减为$1\times10^{-4}$，批量大小为4096，并使用ReduceLROnPlateau调节学习率、梯度范数裁剪5.0和验证集早停。最大训练轮数设置为300轮，模型选择仅依据训练月份内部验证集性能，不使用独立测试月份进行调参。

\subsection{VHP状态监测算例}
\label{subsec:vhp_case_v4}

\subsubsection{研究对象及输入输出配置}

VHP位于主蒸汽进入汽轮机后的首个主要膨胀设备，其入口由A/B侧主蒸汽参数和调节阀状态共同决定，排汽进入一次低温再热系统，同时设置一级抽汽支路向1号高压加热器供汽。因此，VHP既具有明确的主通流边界，又同时包含汽缸本体、抽汽支路及再热接口等多类可观测状态，适合作为设备级健康建模的代表性算例。

实际模型采用16个输入变量。主蒸汽边界包含A/B侧共12个压力和温度测点；操作与工况变量包括2个VHP调节阀实际位置、有功负荷和主蒸汽流量。输入中不加入汽轮机转速，因为正常并网运行阶段转速基本恒定；凝汽器背压属于下游冷端状态，不作为VHP直接输入；高旁状态也不作为本模型输入，只在必要时用于运行工况筛选。

训练阶段设置27个候选状态输出，最终评价与故障诊断使用24点。24点覆盖四类状态区域：VHP本体及主排汽状态、一级抽汽及1号高压加热器入口状态、VHP至一次再热接口状态，以及缸体/内部蒸汽慢变热状态。这样的输出设计使健康模型不仅描述汽缸局部压力温度，还覆盖与后续故障传播相关的直接热力接口。

\subsubsection{数据集及实验设置}

VHP健康状态模型采用2024年5月正常运行数据训练，并以2024年6月数据进行独立测试，采样周期为1 min。经有效值和连续时间窗口筛选后，训练月份可用样本数为43911，独立测试月份可用样本数为42710。为了比较不同动态建模方式，设置ANN、TCN、iTransformer-small和GRU-RNN四种模型，所有模型使用相同的输入输出变量、训练月份和独立测试月份。

ANN采用单时刻输入作为静态基线，不具有显式历史窗口。对于GRU、TCN和iTransformer-small，连续24个采样点构成一个时序输入窗口。仅保留窗口内部相邻时间戳连续的样本，完成序列构造后，GRU训练集包含43759个序列样本，独立测试集包含42595个序列样本。模型输入和输出均采用StandardScaler标准化，所有均值和标准差仅根据训练数据确定。

GRU隐藏维度为32，Dropout率为0.05，训练采用AdamW优化器，初始学习率为$1\times10^{-3}$，权重衰减为$1\times10^{-4}$，batch size为4096，并将梯度范数裁剪至5.0。模型选择阶段根据验证损失使用ReduceLROnPlateau调整学习率，确定最佳训练轮次后，在完整训练月份重新训练最终模型。上述设置用于形成本文稳定基线，不将其声称为经过全局超参数搜索得到的唯一最优组合。

模型采用$R^2$、RMSE和MAE评价单点预测性能：
\begin{equation}
R^2=1-\frac{\sum_{i=1}^{n}(y_i-\hat y_i)^2}{\sum_{i=1}^{n}(y_i-\bar y)^2},
\end{equation}
\begin{equation}
\mathrm{RMSE}=\sqrt{\frac{1}{n}\sum_{i=1}^{n}(y_i-\hat y_i)^2},\qquad
\mathrm{MAE}=\frac{1}{n}\sum_{i=1}^{n}|y_i-\hat y_i|.
\end{equation}
由于24个状态同时包含压力和温度等不同物理量，RMSE和MAE仅用于逐测点或相同量纲参数比较；跨状态总体性能采用无量纲的24点宏平均$R^2$评价。

\subsubsection{模型训练过程}

模型选择阶段训练损失和验证损失随训练轮次整体下降，并在后期进入稳定区间。最低验证MSE为0.057533，出现在epoch 161，对应训练MSE为0.043519。根据该结果，将161轮确定为最终训练轮次；随后利用完整2024年5月训练数据重新训练模型，第161轮完整训练集MSE为0.039283。图2-5同时给出训练和验证损失，并标记最佳验证轮次，以避免仅凭单一训练损失判断模型是否收敛。

训练损失仅反映模型对训练数据的拟合程度，不能单独作为状态监测模型有效性的依据。因此，本文进一步使用完全独立的2024年6月数据评价跨月份泛化性能，重点检验模型在新的运行时间区间和负荷分布下能否继续给出稳定的正常状态参考。

\subsubsection{总体状态监测性能}

在最终24个状态参数上，ANN、TCN、iTransformer-small和GRU-RNN的测试宏平均$R^2$分别为0.9087、0.9397、0.9356和0.9511。GRU取得最高总体精度，相比ANN提高0.0424，相比TCN和iTransformer-small分别提高0.0114和0.0155。需要指出，ANN作为静态基线不具有与三种时序模型相同的历史信息，因此GRU相对ANN的提升同时包含“引入时间历史”和“采用门控时序结构”两方面影响；GRU与TCN、iTransformer-small之间的比较更能反映不同动态特征提取机制的差异。

从状态分组看，GRU在VHP本体及主排汽状态Y1、一级抽汽接口Y2和VHP至RH1接口Y3上的平均$R^2$分别为0.9581、0.9634和0.9569，在缸体与内部蒸汽状态Y4上的平均$R^2$为0.8892。前三组主要由压力、主流蒸汽温度及抽汽/冷再接口状态构成，与入口边界和配汽状态之间具有较直接动态关系；Y4包含金属壁温等慢变状态，受蓄热和更长历史工况影响，其预测难度明显更高。这一分组差异与前述汽轮机参数快慢动态特性分析相一致。

逐测点统计中，GRU在24个状态点中的17个点取得四种模型最高$R^2$，其余7个点由iTransformer-small取得最优结果。由于其中部分压力点各模型本身已达到接近0.998的高精度，单纯统计“最优点数量”并不足以评价模型优劣，因此本文进一步选取具有不同动态特性的代表性状态分析预测曲线。

\subsubsection{典型状态参数预测结果}

对于超高压缸叶片级压力\#1，GRU完整测试月$R^2$达到0.9979，说明主蒸汽边界、阀门状态及负荷信息能够较充分描述该压力状态变化。压力是通流系统中响应较快的状态量，其与入口流量和级组压力分布之间关系相对稳定，因此四种模型在此类测点上通常均能获得较高精度。

对于具有更明显热惯性的排汽蒸汽温度，模型差异更加突出。超高压排汽蒸汽温度\#1上，ANN、TCN、iTransformer-small和GRU的$R^2$分别为0.8092、0.8492、0.8488和0.8676；超高压排汽蒸汽温度\#3上，四种模型对应$R^2$分别为0.8840、0.9153、0.9135和0.9387。与静态ANN相比，三种时序模型均表现更好，说明排汽温度不仅与当前边界有关，还具有明显的时间依赖；GRU在两个代表温度点上进一步取得最高精度，表明门控历史信息对该类动态热状态具有较好的适应性。

一级抽汽接口同样体现出时间特征的作用。一级抽汽逆止门后水平管管顶温度的GRU测试$R^2$为0.9529，高于ANN的0.9056、TCN的0.9385和iTransformer-small的0.9210。该状态位于VHP与回热系统之间，其变化不仅受本体通流状态影响，还与抽汽管路热过程相关；较高预测精度说明设备级模型能够同时建立VHP内部状态和外部热力接口的正常参考。

对于超高压内缸壁温90\%处\#1，GRU测试$R^2$为0.7346，iTransformer-small为0.7483，明显低于压力和主流蒸汽温度状态。金属壁温受到较长时间尺度蓄热过程和历史负荷轨迹影响，24点短时间窗口及当前边界变量尚不能完全描述其长期热状态演化。本文保留该困难测点而不从结果中删除，一方面避免选择性展示造成对模型性能的过度评价，另一方面也明确了当前健康状态模型的主要受限环节：对于高热惯性金属状态，后续可通过扩展历史窗口、增加更直接热边界或引入物理热惯性约束进一步改善。

图2-7选取压力、排汽温度、抽汽接口温度和缸壁温四类代表状态，在统一测试时段展示DCS实测值和GRU预测值。图示区间固定取独立测试集前48 h，仅用于观察局部动态跟踪过程；文中$R^2$均由完整测试月份计算，不根据图示窗口重新选择或计算。

\subsubsection{模型对比与状态监测结果讨论}

综合24点总体精度、状态分组结果和代表性动态曲线，GRU在当前VHP状态监测任务中表现出较好的综合适用性。其优势并非体现在所有测点均绝对最优，而是在较小模型规模下对本体压力、排汽温度、一级抽汽和冷再接口等不同类型状态均保持较稳定预测能力。尤其对于温度和接口状态，历史信息相较静态映射具有更明显价值，这与汽轮机热过程中的滞后和蓄热特性相符。

同时，结果也表明不同状态变量的可建模性存在明显差异。压力和主流蒸汽状态与直接边界之间关系较强，能够形成高精度健康参考；缸体金属温度等慢变状态仍存在较大误差，说明仅依赖短期数据窗口并不足以描述所有时间尺度。该现象并不影响本文采用状态监测模型作为后续故障诊断健康基准的总体思路，但提示后续在构造故障诊断动态状态信息时，应关注不同节点残差尺度和正常波动水平，而不能简单以统一绝对偏差阈值判断异常。

因此，本文最终选择GRU作为汽轮机设备级状态监测的基础模型。其一，GRU在VHP最终24个评价状态上的宏平均$R^2$为0.9511，整体状态覆盖较好；其二，对排汽温度和抽汽接口等具有明显动态特征的状态表现优于静态ANN，并在主要时序对照中保持较好综合性能；其三，网络仅5755个可训练参数，结构紧凑，便于后续在多个设备对象上采用相同建模范式进行扩展。第二章建立的模型输出将作为设备健康运行状态参考，其与实际运行数据之间的偏差信息将在后续KG-GCN故障诊断模型输入构建中进一步处理。

\subsection{本章小结}

本章提出了面向燃煤机组汽轮机系统的设备级健康状态监测方法。首先依据主通流和直接热力接口将系统划分为VHP、再热系统、HP、IP、LP及凝汽器等建模对象，并统一采用$U+B\rightarrow Y$的条件健康映射，避免将上游模型预测继续串联至下游模型。随后介绍GRU动态建模原理、标准化多输出回归、LayerNorm、Dropout和AdamW训练策略，并构建了单层GRU(32)的小型设备模型。

以VHP为代表性算例，采用2024年5月正常DCS数据训练模型，并在完全独立的2024年6月数据上进行测试。最终24个状态参数的测试宏平均$R^2$达到0.9511，GRU在17个状态点上取得四种模型最高$R^2$；对超高压缸叶片级压力、排汽温度和一级抽汽接口温度均形成较稳定健康参考，而金属壁温等慢变状态仍是当前模型的主要困难点。上述结果表明，设备级GRU模型能够在给定运行边界下提供动态正常状态参考，为后续根据实测状态与健康预测之间的偏差信息构造故障诊断动态特征提供基础。